\subsection{Actual polynomial projective ideal sheaves and native quotient points}
\label{subsec:native-homogeneous-quotients}

\begin{theorem}[Actual coordinate section ideal of an original homogeneous quotient]
\label{graded:native-polynomial-quotient-coordinate-ideal}
\lean{EnvelopeProperty.ProjBerk.homogeneousQuotient_coordinateKernel,
EnvelopeProperty.ProjBerk.homogeneousQuotient_coordinateSectionIdeal}
\leanok
\uses{graded:projective-homogeneous-quotient-ideal-sheaf,projective:native-homogeneous-chart-coordinates,projective:native-scheme-coordinate-sections}
Let $R$ be any commutative ring and let $J$ be an actual homogeneous
ideal of $R[X_0,\ldots,X_n]$. On the actual coordinate open $D_+(X_j)$
of polynomial projective space, pull the genuine kernel ideal sheaf of
$\operatorname{Proj}(R[X]/J)$ back under the actual coordinate
section-ring isomorphism. The entire resulting ideal is exactly the
original affine dehomogenization image of $J$, obtained by setting
$X_j=1$ and the other variables to the true coordinate variables.
No radical, reducedness, domain or noetherian hypothesis is required.
\end{theorem}
\begin{proof}
First pull the genuine section ideal back to the actual homogeneous
localization using the preceding whole away-to-section kernel identity.
The actual coordinate ring isomorphism identifies its elements with
these real homogeneous fractions, not merely their coordinate values.

Suppose an actual coordinate element lies in that kernel. Express its
homogeneous-localization image as a same-degree fraction $F/X_j^d$.
The genuine original-ideal annihilator criterion gives $X_j^mF\in J$
for some $m$. Under the actual dehomogenization ring map, $X_j$ becomes
one, so the entire polynomial $F$ dehomogenizes to an element of the
actual image ideal of $J$. The true inverse coordinate-ring identity
identifies that polynomial with the original selected coordinate
element. This proves the first ideal inclusion.

Conversely, for $F\in J$ decompose the entire polynomial into its
actual homogeneous components. Each lies in the original ideal by
its genuine homogeneity. The actual dehomogenized value of each
component corresponds to the real homogeneous fraction with that
component as numerator. Its quotient image is zero by the original
annihilator criterion with exponent zero. Add the entire finite
homogeneous decomposition and use the kernel ideal's closure under
multiplication to contain all of the actual image ideal. Thus both
whole ideal inclusions hold. The true homogeneous-fraction section
identity then gives the stated actual section-ideal comparison.
\end{proof}

\begin{theorem}[Actual native homogeneous quotient and original normalized vanishing locus]
\label{graded:native-polynomial-homogeneous-quotient-points}
\lean{EnvelopeProperty.ProjBerk.homogeneousQuotient_normalizedClosedLocus,
EnvelopeProperty.ProjBerk.nativeHomogeneousQuotientHomeomorph,
EnvelopeProperty.ProjBerk.nativeHomogeneousQuotientHomeomorph_apply,
EnvelopeProperty.ProjBerk.nativeHomogeneousQuotientHomeomorph_support,
EnvelopeProperty.ProjBerk.compactSpace_nativeHomogeneousQuotient,
EnvelopeProperty.ProjBerk.t2Space_nativeHomogeneousQuotient}
\leanok
\uses{graded:projective-homogeneous-quotient-ideal-sheaf,projective:native-closed-normalized-comparison,projective:native-closed-projective-compactness}
For any prescribed real seminorm $v$ on $R$, the actual native point
space of $\operatorname{Proj}(R[X]/J)$ above $v$ is homeomorphic to
the normalized graded seminorms of $R[X]$ with the original coefficient
values $v$ and with $J$ contained in their entire polynomial kernel.
The comparison preserves the whole actual native inclusion and the
entire homogeneous prime support. The actual native quotient space
is compact and Hausdorff for every $v$, including unbounded coefficient
seminorms; the locus can be empty. No field, completion or nonempty
quotient hypothesis is required.
\end{theorem}
\begin{proof}
Use the actual homogeneous quotient closed immersion proved above and
its genuine source structure map obtained by composition with the
original polynomial projective base morphism. The independently
constructed native closed-projective homeomorphism identifies its
entire fixed-base source point space with the normalized points whose
actual native support lies in its true scheme image. The proved
original-ideal range equality makes this condition exactly containment
of $J$ in the entire polynomial seminorm kernel; this support equality
holds on whole ideals, without a radical or chart certificate.
Restrict the genuine homeomorphism by this literal set equality.
Its whole-point formula and polynomial support formula are the same
as those of the underlying actual closed-projective comparison.
Finally its source compactness and Hausdorffness are the actual
closed-projective producers, which use projective compactness for
every fixed $v$ and the genuine native closed embedding. No uniform
bound on the coefficients is introduced.
\end{proof}
